% SPDX-License-Identifier: GPL-3.0-or-later
% Copyright (C) 2026 Alessandro Greco
%% Testo alternativo per l'accessibilita' (solo EPUB).
\altfig{Schema semplificato del traffico Tor su TCP e TLS. L'osservatore esterno può distinguere indirizzi IP, porte, volumi, temporizzazione e lunghezze dei record TLS, ma senza chiavi TLS non leggere le celle. Il flusso TLS può contenere più celle e una cella può attraversare più record: non esiste corrispondenza obbligatoria fra confini di record e celle. Con il protocollo di collegamento v4 o superiore, le celle di lunghezza fissa sono di 514 byte: intestazione da 5 byte e corpo da 509 byte; esistono anche celle di lunghezza variabile. Solo il corpo delle celle RELAY è cifrato a strati con le chiavi del circuito: PADDING e VPADDING non sono messaggi RELAY. Con le chiavi TLS si possono esaminare le celle, ma non automaticamente decifrare il corpo RELAY.}
% Schema di incapsulamento (i confini dei record e delle celle non coincidono).
\begin{tikzpicture}[font=\sffamily\scriptsize,
  liv/.style={draw,rounded corners=2pt,minimum height=0.85cm,align=center},
  cel/.style={draw,fill=torviola!12,minimum height=0.7cm,minimum width=1.5cm,font=\sffamily\tiny}]
  \node[liv,fill=torgrigio!10,minimum width=10.6cm] (tcp) at (0,2.6) {flusso TCP: segmenti variabili};
  \node[liv,fill=torverde!12,minimum width=9.4cm] (tls) at (0,1.35) {TLS: record non allineati alle celle};
  \foreach \i/\x in {1/-3.6,2/-1.8,3/0,4/1.8,5/3.6} \node[cel] (c\i) at (\x,0) {514\,B};
  \node[anchor=west,font=\sffamily\tiny,text=torgrigio] at (4.6,0) {fisse v4+};
  \node[liv,fill=avviso!10,minimum width=4.6cm] (rel) at (-1.7,-1.4) {RELAY: corpo 509\,B\\cifrato a strati};
  \node[liv,fill=torgrigio!8,minimum width=4.6cm] (pad) at (3.3,-1.4) {PADDING / VPADDING\\non RELAY};
  \draw[torgrigio] (c1.south west) -- (rel.north west); \draw[torgrigio] (c3.south east) -- (rel.north east);
  \node[anchor=west,align=left,font=\sffamily\tiny,text=torgrigio] at (5.8,2.6) {IP, porte, volumi\\e tempi osservabili};
  \node[anchor=west,align=left,font=\sffamily\tiny,text=torgrigio] at (5.8,1.35) {lunghezze dei record\\visibili sul canale};
  \node[anchor=west,align=left,font=\sffamily\tiny,text=torgrigio] at (5.8,0) {celle dopo la\\decifratura TLS};
  \node[anchor=west,align=left,font=\sffamily\tiny,text=torgrigio] at (5.8,-1.4) {RELAY: occorrono\\chiavi del circuito};
  \node[font=\sffamily\tiny,text=torgrigio,align=center] at (0,-2.55) {Schema: i confini TLS non coincidono necessariamente con quelli delle celle.};
\end{tikzpicture}
